% Appendix A. Appendices are ordinary chapters after \appendix in main.tex.
\chapter{Reference test protocol}
\label{app:protocol}

Every reference performance test followed the sequence in
\cref{tab:protocol}. The site was disconnected from the market at least
six hours in advance so that the cells could relax, and the ambient
temperature of the container was held at 25\,°C by its climate control.
Capacity is the charge delivered in step~4, divided by the nameplate
capacity at commissioning.

\begin{table}[htb]
  \caption[Reference test sequence]{Reference test sequence. C/5 is the
    current that would discharge the nameplate capacity in five hours.}
  \label{tab:protocol}
  \centering
  \begin{tabular}{@{}r l l l@{}}
    \toprule
    Step & Action & Current & End condition \\
    \midrule
    1 & Rest            & ---      & 4 h, cell spread \(<\) 1\,°C \\
    2 & Charge (CC--CV) & C/5      & upper voltage, then \(I <\) C/50 \\
    3 & Rest            & ---      & 1 h \\
    4 & Discharge (CC)  & C/5      & lower voltage limit \\
    5 & Rest            & ---      & 1 h \\
    6 & Recharge        & C/5      & 50\,\% state of charge \\
    \bottomrule
  \end{tabular}
\end{table}

Two corrections were applied. Charge was converted to capacity at a
reference temperature of 25\,°C using the manufacturer's temperature
coefficient, and the auxiliary consumption of the container, measured
separately for each site, was subtracted from the pack current. The
repeatability of the test, estimated from back-to-back repetitions on
two sites, is \(\pm 0.4\) percentage points of state of health.
